\documentclass[10pt,a4paper]{article}
\usepackage[utf8]{inputenc}
\usepackage[T1]{fontenc}
\usepackage[french]{babel}
\usepackage[margin=1.7cm]{geometry}
\usepackage{booktabs}
\usepackage{array}
\usepackage{enumitem}
\setlist{nosep, leftmargin=1.3em}
\setlength{\parindent}{0pt}
\setlength{\parskip}{2pt}
\newcommand{\code}[1]{\texttt{\small #1}}
\pagestyle{empty}

\title{\vspace{-1.3cm}\textbf{Scénario d'évaluation : Prototype (migration de connexion) vs Vanilla}\\
\large Bancs d'essai issus de SeBS, pilotés par \code{wrk2} (HTTPS, keep-alive)}
\author{}\date{}

\begin{document}
\maketitle
\vspace{-1.4cm}

\section{Les applications (banc SeBS)}
Les quatre fonctions sont issues du banc de référence \textbf{SeBS} (\emph{Serverless Benchmark
Suite}, SPCL/ETH Zürich, Middleware 2021) : leur \textbf{code métier est repris verbatim}, seul
l'hôte (Flask + gunicorn derrière le watchdog) et le pilotage (notre harnais \code{wrk2}) sont à nous.
Elles couvrent des profils de charge variés et volontairement légers.

\begin{center}\small
\begin{tabular}{@{}llll@{}}
\toprule
\textbf{Application} & \textbf{Réf.\ SeBS} & \textbf{Aspect évalué} & \textbf{Stockage} \\
\midrule
\code{dynamic-html}   & 110.dynamic-html   & CPU / templating web (grosse réponse) & aucun \\
\code{thumbnailer}    & 210.thumbnailer    & CPU image + E/S objet                 & MinIO \\
\code{compression}    & 311.compression    & CPU compression + E/S objet           & MinIO \\
\code{graph-pagerank} & 501.graph-pagerank & CPU pur (igraph, PageRank)            & aucun \\
\bottomrule
\end{tabular}
\end{center}

\section{Entrées et sorties par application}
Chaque requête POSTe l'\emph{event} JSON ci-dessous ; pour \code{thumbnailer} et \code{compression},
les données sont d'abord \textbf{semées dans MinIO} par \code{seed-storage.py} (ici : image JPEG
\textbf{16$\times$12} $\approx$739~o, et dataset de \textbf{12 fichiers} $\approx$48~KB). La sortie est
celle du \textbf{code SeBS d'origine}.

\begin{center}\footnotesize
\begin{tabular}{@{}l p{6.9cm} p{6.6cm}@{}}
\toprule
\textbf{App.} & \textbf{Entrée (event + données semées)} & \textbf{Sortie (SeBS)} \\
\midrule
\code{dynamic-html} &
\code{username} : nom inséré dans la page ; \code{random\_len=1000} : nb d'entiers aléatoires
(0--10$^6$) générés $\rightarrow$ pilote la \textbf{taille du HTML} et la charge CPU. Aucune donnée externe. &
Page \textbf{HTML} rendue via Jinja2 (username + heure + liste des nombres). \code{\{result: html\}}. \\
\addlinespace
\code{graph-pagerank} &
\code{size=200} : nb de n\oe{}uds du graphe Barab\'asi--Albert généré ; \code{seed=42} : graine
aléatoire (reproductibilité). &
Score \textbf{PageRank} du 1\textsuperscript{er} n\oe{}ud + \code{measurement}
(\code{graph\_generating\_time}, \code{compute\_time}). \\
\addlinespace
\code{thumbnailer} &
\code{bucket}/\code{input}/\code{output} : bucket et préfixes MinIO ; \code{object.key=sample.jpg} :
image source ; \code{width=10, height=10} : \textbf{taille cible de la miniature}. &
\textbf{Miniature JPEG} redimensionnée, ré-uploadée dans \code{thumbnailer/output/} ;
\code{result=\{bucket,key\}} + \code{measurement} (temps/tailles download, compute, upload). \\
\addlinespace
\code{compression} &
\code{bucket}/\code{input}/\code{output} : bucket et préfixes ; \code{object.key=dataset} :
\textbf{sous-dossier à compresser}. &
\textbf{Archive ZIP} du dataset, ré-uploadée dans \code{compression/output/} ;
\code{result=\{bucket,key\}} + \code{measurement} (idem). \\
\bottomrule
\end{tabular}
\end{center}

\section{Scénario d'évaluation}
Pour chaque application, on effectue un \textbf{balayage de débit} (\emph{rate sweep}) en boucle
ouverte : \code{wrk2} envoie l'\emph{event} JSON de la fonction (POST, keep-alive) à un \textbf{débit
cible constant}, palier par palier. Chaque palier dure \textbf{60~s}, suivi d'une \textbf{pause de 5~s}
pour laisser le système revenir au repos. À chaque palier on relève : RPS atteint, latences
(moy/p50/p99), CPU et réseau du Pi, coût du chemin de données (\emph{client/server/overhead}) et les
erreurs. La \textbf{même charge exacte} est rejouée dans les deux modes ; seul le mode change :
\begin{itemize}
\item \textbf{Prototype} : le gateway \textbf{migre la connexion} (descripteur TCP + état TLS 1.3) vers
le conteneur, qui répond \textbf{directement} au client (\code{--mode proto}).
\item \textbf{Vanilla} : chaîne proxy OpenFaaS standard, le gateway relaie requête et réponse
(\code{--mode vanilla}).
\end{itemize}
Les images des fonctions sont \textbf{identiques} dans les deux cas ; le comportement du gateway est
fixé par sa configuration (\code{HTTPMIGRATE\_ENABLE}), et \code{--mode} n'est qu'une étiquette côté
client. Toutes les mesures sont en \textbf{HTTPS} (TLS~1.3, wolfSSL, certificat auto-signé) sur le
port~8443.

\section{Paramètres \code{wrk2} par application}
Paramètres \textbf{communs} à toutes les applications et aux deux modes : \code{--scheme https}
(port~8443), \code{--host 192.168.2.2} (le Pi), \code{--duration-s 60}, \code{--timeout-s 60},
\code{--pause 5}. Le client wrk2 ouvre \emph{concurrency} connexions réparties sur \emph{threads}
threads, et vise le débit \code{-R} de chaque palier.

\begin{center}\small
\begin{tabular}{@{}lp{6.6cm}cc@{}}
\toprule
\textbf{Application} & \textbf{Débits ciblés (\code{--rates}, req/s)} & \textbf{Concurrence (\code{-c})} & \textbf{Threads (\code{-t})} \\
\midrule
\code{compression}    & 2, 4, 6, 8, 10, 12, 14, 16                                   & 16 & 4 \\
\code{graph-pagerank} & 32, 48, 64, 80, 96, 112, 128, 144, 160, 176, 192, 208, 224, 240, 256, 272, 288 & 80 & 8 \\
\code{thumbnailer}    & 1, 16, 32, 48, 64, 80, 96, 112, 128, 144, 160, 176, 192      & 32 & 4 \\
\code{dynamic-html}   & 1, 16, 32, 48, 64, 80, 96, 112, 128, 144, 160, 176, 192      & 32 & 4 \\
\bottomrule
\end{tabular}
\end{center}
Les plages diffèrent car les applications saturent à des débits très différents : \code{compression}
et \code{graph-pagerank} sont plus lourdes par requête (débits utiles plus bas et concurrence adaptée),
tandis que \code{dynamic-html} sature le CPU à haut débit.

\section{Environnement d'évaluation}
\begin{itemize}
\item \textbf{Serveur} : Raspberry~Pi~4 (\code{192.168.2.2}), 4 c\oe{}urs ARM Cortex-A72, carte réseau
$\approx$~110~Mo/s (940~Mbit/s mesurés à \code{iperf3}). Il exécute \textbf{faasd} + le gateway OpenFaaS
modifié (terminaison/migration TLS via wolfSSL).
\item \textbf{Fonctions} : hôte Python (Flask + \textbf{gunicorn, 4 workers} $\rightarrow$ vrai
parallélisme multi-c\oe{}ur) derrière le watchdog \emph{full-proxy}. Stockage objet \textbf{MinIO}
(\code{10.62.0.1:9000}) pour \code{thumbnailer} et \code{compression}.
\item \textbf{Client} : machine de charge exécutant \code{wrk2} (générateur en boucle ouverte, correction
de \emph{coordinated omission}), POST de l'event JSON à débit constant, connexions keep-alive.
\item \textbf{Instrumentation} (par palier) : latences via wrk2 ; \textbf{CPU} du Pi via \code{sar -u}
(échelle 0--100\,\%) et \textbf{réseau} via \code{sar -n DEV} sur \code{eth0} ; \emph{perf-cost}
(\code{client\_ms}, \code{server\_ms} du wrapper SeBS, \code{overhead = client - server}) ; comptage des
réponses et des erreurs/timeouts.
\end{itemize}

\end{document}
